%Autor: miguev

\pagestyle{empty}

\vspace*{20mm}
\begin{center}
{\LARGE Gu�a de aprendizaje}\vspace*{5mm}
%\resizebox{\textwidth}{10mm}{{\LARGE\bf Curso de Introducci�n a Linux para Alumnos}}
\resizebox{\textwidth}{10mm}{{\LARGE\bf Software Libre para estudiantes}}
\end{center}

\vfill
\cleardoublepage

\vspace*{20mm}
\begin{center}
{\LARGE Gu�a de aprendizaje}\vspace*{5mm}
%\resizebox{\textwidth}{10mm}{{\LARGE\bf Curso de Introducci�n a Linux para Alumnos}}
\resizebox{\textwidth}{10mm}{{\LARGE\bf Software Libre para estudiantes}}
\end{center}

\begin{center}
\vspace*{60mm}
\includegraphics[width=120mm]{imagenes/gulic.eps}
\end{center}

\vfill

\begin{flushright}
%{\LARGE Grupo de Usuarios de Linux de Canarias}
%\vspace*{5mm}
{\large
Miguel �ngel Vilela Garc�a     \\ % 69 pags
Jes�s Miguel Torres Jorge      \\ % 65 pags
Carlos P�rez P�rez             \\ % 51 pags
Tom�s Bautista Delgado         \\ % 50 pags
Carlos Alberto Morales D�az    \\ % 44 pags
F�lix J. Marcelo Wirnitzer     \\ % 30 pags
Carlos de la Cruz Pinto        \\ % 18 pags
Ren� Mart�n Rodr�guez          \\ % 16 pags
Ed�n Kozo                      \\ % 13 pags
Carlos Mestre Gonz�lez         \\ % 11 pags
% Teresa Gonz�lez de la F�       \\ %  9 pags
% Luis Cabrera Sa�co             \\ %  8 pags
% Pedro Gracia Fajardo           \\ %  2 pags
% Carlos A. Paramio              \\ %  1 pags
}
\end{flushright}

% Autores ordenados por cantidad de p�ginas:
% Miguel �ngel Vilela Garc�a
% Jes�s Miguel Torres Jorge
% Carlos P�rez P�rez
% Tom�s Bautista Delgado
% Carlos Alberto Morales D�az
% F�lix J. Marcelo Wirnitzer
% Carlos de la Cruz Pinto
% Ren� Mart�n Rodr�guez
% Ed�n Kozo
% Carlos Mestre Gonz�lez
% Teresa Gonz�lez de la F�
% Luis Cabrera Sa�co
% Pedro Gracia Fajardo
% Carlos A. Paramio

\vfill
\clearpage

\vspace*{130mm}
\noindent Copyright  \copyright~ 2001 -- 2003 \GULiC
\vspace*{5mm}

\noindent  Se da  permiso  para copiar,  distribuir  o modificar  este
documento  en los  t�rminos que  establece la  GNU Free  Documentation
License, Versi�n 1.1 o cualquier  otra versi�n posterior publicada por
la Free Software  Foundation; las secciones que  no pueden modificarse
son  ``Presentaci�n''  y  ``Agradecimientos'',  y  no  existen  textos
cubierta ni contracubierta. Se incluye una  copia de la licencia en la
secci�n ``GNU Free Documentation  License''. Este libro es \copyright~
de sus autores.

\noindent  La forma  original de  este documento  es c�digo  fuente en
\LaTeX.  La compilaci�n  de este  c�digo  fuente en  \LaTeX~ tiene  el
efecto de  generar una representaci�n del  documento independiente del
dispositivo, que puede convertirse a otros formatos o imprimirse.

\noindent  La  GNU  Free  Documentation  License  est�  disponible  en
www.gnu.org o solici�ndola por escrito a the Free Software Foundation,
Inc., 59 Temple Place - Suite 330, Boston, MA 02111-1307, USA.

\noindent  Este  documento lo  han  maquetado  sus autores  utilizando
\LaTeX~ y {\sf The GIMP}, que son programas abiertos y libres.

\vfill
\clearpage

\vspace*{35mm}
\begin{flushright}
{\large A los novatos y novatas en el mundo de GNU/Linux, \\
y a sus familiares, amistades y parejas} \\
{\small (ellas y ellos saben por qu�)}
\end{flushright}

\vfill
\cleardoublepage

\pagestyle{fancy}
